\documentclass[12pt,a4paper]{article}
\usepackage[utf8]{inputenc}
\usepackage[T1]{fontenc}
\usepackage[spanish,es-nodecimaldot]{babel}
\usepackage[a4paper,margin=2.45cm,headheight=26pt]{geometry}
\usepackage{lmodern}
\usepackage{graphicx}
\usepackage{xcolor}
\usepackage{booktabs}
\usepackage{longtable}
\usepackage{listings}
\usepackage{tikz}
\usetikzlibrary{positioning,arrows.meta,calc}
\usepackage{float}
\usepackage{caption}
\usepackage{tcolorbox}
\usepackage{fancyhdr}
\usepackage{enumitem}
\usepackage{hyperref}

\definecolor{techpurple}{HTML}{7135A5}
\definecolor{techpurple2}{HTML}{914BC6}
\definecolor{techlavender}{HTML}{F3EFF8}
\definecolor{techink}{HTML}{25212B}
\definecolor{techpink}{HTML}{D94E91}
\definecolor{techline}{HTML}{DED0EA}
\hypersetup{colorlinks=true,linkcolor=techpurple,urlcolor=techpurple2,citecolor=techpurple}
\pagestyle{fancy}
\fancyhf{}
\lhead{\textcolor{techpurple}{\textbf{TECH-FLOW}}}
\rhead{\textcolor{techpurple}{\leftmark}}
\cfoot{\textcolor{techpurple}{\thepage}}
\renewcommand{\headrulewidth}{0.5pt}
\renewcommand{\headrule}{\hbox to\headwidth{\color{techline}\leaders\hrule height \headrulewidth\hfill}}
\setlength{\parskip}{5pt}
\setlength{\parindent}{0pt}
\setlength{\emergencystretch}{3em}
\setlist{itemsep=3pt,topsep=4pt}
\captionsetup{font=small,labelfont={bf,color=techpurple},labelsep=colon}
\lstset{basicstyle=\ttfamily\small,breaklines=true,frame=single,rulecolor=\color{techline},backgroundcolor=\color{techlavender},showstringspaces=false,keywordstyle=\color{techpurple}\bfseries,commentstyle=\color{gray}}
\newtcolorbox{techbox}[1][]{colback=techlavender,colframe=techline,arc=3mm,boxrule=.6pt,left=7pt,right=7pt,top=6pt,bottom=6pt,#1}
\newtcolorbox{warningbox}{colback=techpink!7,colframe=techpink,arc=3mm,boxrule=.6pt,left=7pt,right=7pt,top=6pt,bottom=6pt}
\newcommand{\TechFigure}[4]{%
\begin{figure}[H]\centering
\includegraphics[width=\linewidth,height=#2,keepaspectratio]{#1}
\caption{#3}\label{#4}
\end{figure}}
\newcommand{\TechCover}[3]{%
\begin{titlepage}
\begin{tikzpicture}[remember picture,overlay]
\fill[techlavender] (current page.south west) rectangle (current page.north east);
\fill[techpurple] (current page.north west) rectangle ([yshift=-3.8cm]current page.north east);
\fill[techpurple2] ([xshift=-3.4cm]current page.south east) circle (4.4cm);
\draw[techpurple2,line width=2pt] ([xshift=1.5cm,yshift=-7.4cm]current page.north west) circle (2.2cm);
\draw[white,line width=1.5pt] ([xshift=1.5cm,yshift=-7.4cm]current page.north west) circle (1.45cm);
\end{tikzpicture}
\vspace*{0.65cm}\color{white}\raggedright
{\fontsize{24}{28}\selectfont\bfseries TECH-FLOW\par}
{\large Sistema de gestión de inventario y ventas\par}
\vspace{3.2cm}\color{techink}
{\fontsize{30}{35}\selectfont\bfseries #1\par}
\vspace{0.35cm}{\Large\color{techpurple}#2\par}
\vfill
\begin{techbox}[colback=white,colframe=techline]
\textbf{Universidad:} Universidad Mayor de San Simón\\
\textbf{Docente:} CORINA JUSTINA FLORES VILLARROEL\\
\textbf{Integrantes:}\\
Jorge Henry Quispe Escalera\\
Logan Betho Canelas Silvestre\\
Fausto Jafeth Vilches Mendieta\\
Eunice Melany Zurita Pérez\\
\textbf{Entrega:} 30 de noviembre de 2026\\
\textbf{Defensas:} 1, 2 y 8 de diciembre de 2026\\
\end{techbox}
\end{titlepage}}

\begin{document}
\TechCover{MANUAL DE INSTALACIÓN DEMO}{Preparación, compilación y ejecución}{MANUAL DE INSTALACIÓN DEMO}
\tableofcontents\newpage
\section{Antes de comenzar}
Este manual permite preparar Tech-Flow desde una copia limpia del proyecto. Los comandos se han contrastado con el \texttt{pom.xml}, que declara Java 17, JavaFX 17.0.2, JUnit 5.10.0 y el plugin de ejecución JavaFX.
\begin{techbox}\textbf{Ruta de trabajo.} Todos los comandos deben ejecutarse desde la carpeta raíz que contiene \texttt{pom.xml}, \texttt{src} y \texttt{data}.\end{techbox}
\section{Requisitos técnicos}
\begin{table}[H]\centering\caption{Requisitos comprobados}\begin{tabular}{p{.3\textwidth}p{.55\textwidth}}\toprule Componente & Requisito\\\midrule
Sistema operativo & Windows, macOS o Linux con JDK compatible\\
Java & JDK 17\\Maven & Apache Maven 3.8 o superior recomendado\\
Dependencias & Internet durante la primera compilación\\
Espacio & Aproximadamente 1 GB libre\\\bottomrule\end{tabular}\end{table}
\section{1. Obtener el proyecto}
Copie o clone la carpeta \texttt{Tech-Flow}. En Windows PowerShell:
\begin{lstlisting}
cd C:\ruta\hasta\trabajo\Tech-Flow
\end{lstlisting}
Confirme que existen \texttt{pom.xml}, \texttt{src} y \texttt{data}.
\section{2. Verificar Java y Maven}
\begin{lstlisting}
java -version
mvn -version
\end{lstlisting}
Java debe indicar versión 17. Si Maven no encuentra Java, configure \texttt{JAVA\_HOME} al directorio del JDK.
\section{3. Compilar y ejecutar pruebas}
\begin{lstlisting}
mvn clean test
\end{lstlisting}
Los reportes JUnit aparecen en \texttt{target/surefire-reports}.
\begin{warningbox}\textbf{Importante.} No borre ni sustituya los CSV de \texttt{data/} sin realizar una copia de seguridad.\end{warningbox}
\section{4. Ejecutar Tech-Flow}
\begin{lstlisting}
mvn javafx:run
\end{lstlisting}
El plugin utiliza \texttt{com.overbyte.Main}. La ejecución desde la raíz permite localizar los recursos FXML y los datos.
\section{5. Verificar la instalación}
Debe aparecer la pantalla de inicio de sesión. Compruebe que los botones de rol y el campo de contraseña responden. Las credenciales deben corresponder a los registros disponibles en \texttt{data/empleado.csv}.
\TechFigure{../capturas/01-login.png}{7.3cm}{Pantalla que confirma que la aplicación JavaFX inició correctamente.}{fig:instalacion-login}
\section{Errores frecuentes}
\begin{description}\item[Java o Maven no encontrados] Instale un JDK 17 y Maven y añada sus ejecutables al PATH.\item[Dependencias fallan] Repita la compilación con Internet disponible.\item[No aparecen datos] Ejecute desde la raíz y confirme \texttt{data/}.\item[JavaFX no inicia] Use \texttt{mvn javafx:run}, no un classpath manual.\end{description}
\section{Desinstalación}
Cierre la aplicación y elimine la carpeta del proyecto. La caché local de Maven es opcional: borrarla puede afectar a otros proyectos.
\section{Evidencias no proporcionadas}
Se dispone de una captura real del inicio de la aplicación, incluida en la sección \ref{fig:instalacion-login}. No se proporcionaron capturas de la instalación de Java, la terminal con \texttt{mvn clean test} ni la ejecución de \texttt{mvn javafx:run}; esas evidencias deben capturarse en el entorno de instalación utilizado si la universidad las exige.
\end{document}
